\paragraph{Endpoint Tristate Compression Principle: \texorpdfstring{$\mathcal S_{\star}=\{0,2,-2\}$}{S*={0,2,-2}} and the Modulo \texorpdfstring{$4$}{4} Three-Channel System}
To further compress the \(\ZZ/4\) residue from Proposition \ref{prop:half-angle-z4-residue} into a "finite-state" template, we introduce the endpoint tristate set.

\begin{definition}[Endpoint Tristate Set and Endpoint State Function]\label{def:endpoint-tristate-set}
At the endpoint \(w=-1\) (equivalently \(u=-1\) and \(S=0\)), denote
\[
\boxed{\ \mathcal S_{\star}:=\{0,2,-2\}\ }.
\]
The mapping
\[
d\ \longmapsto\ C_d(0)\in\mathcal S_{\star}\qquad(d\ge 1)
\]
is called the \textbf{endpoint state function} of the power-dilation exponent \(d\).
\end{definition}

\begin{corollary}[Endpoint Tristate Compression Principle (\texorpdfstring{$\bmod 4$}{mod 4} Three-Channel Degeneration)]\label{cor:endpoint-tristate-compression}
For all integers \(d\ge 1\), the endpoint values satisfy
\begin{equation}\label{eq:endpoint-tristate-values}
\boxed{\;
C_d(0)=\mathrm{i}^{d}+\mathrm{i}^{-d}=2\cos(d\pi/2)\in\mathcal S_{\star},\qquad
C_d(0)=
\begin{cases}
0, & d\ \text{is odd},\\
2, & d\equiv 0\ (\mathrm{mod}\ 4),\\
-2,& d\equiv 2\ (\mathrm{mod}\ 4).
\end{cases}\;}
\end{equation}
Therefore, the completion dynamics \(S\mapsto C_d(S)\) induced by any power substitution semigroup \(\{u\mapsto u^d\}_{d\ge 1}\) necessarily degenerates on the \(0\)-jet at the endpoint \(S=0\) into a modulo \(4\) tristate system.
\end{corollary}
\begin{proof}
Equation \eqref{eq:endpoint-tristate-values} is simply the expansion of the endpoint identity \eqref{eq:chebyshev-endpoint-values} from Proposition \ref{prop:half-angle-z4-residue} under the three residue classes \(d\bmod 4\). \qed
\end{proof}

\paragraph{Witt--Chebyshev Endpoint Jet Decomposition: Tristate Template of the Primitive}
Let \(p_n(u)\in\ZZ[u]\) be the primitive polynomial, and define in the completion coordinates \(u=r^2,\ S=r+r^{-1}\)
\(\widehat p_n(S):=r^{-n}p_n(r^2)\in\ZZ[S]\);
simultaneously, for any \(d\ge 1\) denote \(C_d(S):=r^d+r^{-d}\in\ZZ[S]\).
Then the strict doubling-angle law \(r^{-dn}p_n(r^{2d})=\widehat p_n(C_d(S))\) holds (see Corollary \ref{cor:primitive-completion-hatp}).
Accordingly, we introduce the endpoint tristate values
\begin{equation}\label{eq:primitive-endpoint-tristate-values}
\boxed{\;
A_n:=\widehat p_n(0),\qquad
B_n^{+}:=\widehat p_n(2),\qquad
B_n^{-}:=\widehat p_n(-2)\; }.
\end{equation}

\begin{proposition}[Tristate Degeneration of the Primitive: \texorpdfstring{$\widehat p_n(C_d(0))$}{hat p_n(C_d(0))} Modulo \texorpdfstring{$4$}{4} Channel Splitting]\label{prop:primitive-endpoint-tristate-degeneration}
For any \(n\ge 1\) and \(d\ge 1\), at the endpoint \(S=0\) the strict degeneration
\begin{equation}\label{eq:primitive-endpoint-tristate-degeneration}
\boxed{\;
\widehat p_n\!\bigl(C_d(0)\bigr)=
\begin{cases}
A_n, & d \text{ odd},\\
B_n^{+}, & d\equiv 0\ (\mathrm{mod}\ 4),\\
B_n^{-}, & d\equiv 2\ (\mathrm{mod}\ 4).
\end{cases}\;}
\end{equation}
holds. Moreover, the identification between the endpoint values \(\pm 2\) and \(u=1\) yields
\begin{equation}\label{eq:primitive-endpoint-Bpm-identification}
\boxed{\ B_n^{+}=p_n(1),\qquad B_n^{-}=(-1)^n p_n(1)\ }.
\end{equation}
\end{proposition}
\begin{proof}
By Corollary \ref{cor:endpoint-tristate-compression}, the endpoint state \(C_d(0)\) takes only the three values \(\{0,2,-2\}\), determined by \(d\bmod 4\); substituting into \(\widehat p_n\) immediately yields \eqref{eq:primitive-endpoint-tristate-degeneration}.
\par
On the other hand, by definition \(\widehat p_n(S)=r^{-n}p_n(r^2)\). Setting \(S=2\) corresponds to \(r=1\), giving \(\widehat p_n(2)=p_n(1)\).
Setting \(S=-2\) corresponds to \(r=-1\), giving \(\widehat p_n(-2)=(-1)^{-n}p_n(1)=(-1)^n p_n(1)\).
This yields \eqref{eq:primitive-endpoint-Bpm-identification}. \qed
\end{proof}

\paragraph{Endpoint Witt Decomposition Template: Three-Channel Weight Kernel and Logarithmic Principal Scale}
For a function \(g\) analytic on some open neighborhood containing the interval \([-2,2]\), define the truncated completion power chain operator
\begin{equation}\label{eq:completed-witt-pow-operator-trunc}
\boxed{\;
(W_{\mathrm{pow}}^{\le M}g)(S):=\sum_{d=1}^{M}\frac{1}{d}\,g\!\bigl(C_d(S)\bigr)\qquad(M\ge 1)\; }.
\end{equation}
Then the \(0\)-jet at the endpoint \(S=0\) satisfies a strict modular channel decomposition:

\begin{proposition}[Endpoint Three-Channel Decomposition: Truncated Power Chain at \texorpdfstring{$S=0$}{S=0} Modulo \texorpdfstring{$4$}{4} Grouping Identity]\label{prop:endpoint-mod4-channel-decomposition}
Let
\[
H_M^{(1)}:=\sum_{\substack{1\le d\le M\\ d\ \mathrm{odd}}}\frac1d,\qquad
H_M^{(0)}:=\sum_{\substack{1\le d\le M\\ d\equiv 0\ (\mathrm{mod}\ 4)}}\frac1d,\qquad
H_M^{(2)}:=\sum_{\substack{1\le d\le M\\ d\equiv 2\ (\mathrm{mod}\ 4)}}\frac1d.
\]
Then for any analytic \(g\) the identity
\begin{equation}\label{eq:endpoint-mod4-channel-decomposition}
\boxed{\;
(W_{\mathrm{pow}}^{\le M}g)(0)
=
H_M^{(1)}\,g(0)
+H_M^{(0)}\,g(2)
+H_M^{(2)}\,g(-2)\; }.
\end{equation}
holds. Therefore, at the endpoint the power chain superposition degenerates at the \(0\)-jet level into three mutually decoupled modular channels, with weights entirely determined by the kernels \(\sum_{d\equiv a\ (\mathrm{mod}\ 4)}\frac1d\).
\end{proposition}
\begin{proof}
By definition \eqref{eq:completed-witt-pow-operator-trunc} and the endpoint tristate compression \eqref{eq:endpoint-tristate-values},
classifying each term \(g(C_d(0))\) according to \(d\bmod 4\) and combining yields \eqref{eq:endpoint-mod4-channel-decomposition}. \qed
\end{proof}

\begin{corollary}[Structural Vanishing of the Odd Channel (Endpoint \(0\)-Jet)]\label{cor:endpoint-odd-channel-vanish}
If the analytic observable \(g\) satisfies \(g(0)=0\), then for any \(M\ge 1\),
\[
\boxed{\;
(W_{\mathrm{pow}}^{\le M}g)(0)=H_M^{(0)}\,g(2)+H_M^{(2)}\,g(-2)\; }.
\]
Therefore, the \(g(0)\) component carried by the odd-\(d\) channel among the endpoint three channels vanishes completely; the \(0\)-jet information of the endpoint boundary layer is provided solely by the comparison of the two even channels \(S=\pm 2\).
\end{corollary}
\begin{proof}
Substituting \(g(0)=0\) into \eqref{eq:endpoint-mod4-channel-decomposition} yields the result immediately. \qed
\end{proof}

\begin{proposition}[Dirichlet-Regularized Weight Kernels and Closed-Form Representation of the Endpoint Three Channels]\label{prop:endpoint-dirichlet-symbol-weights}
For \(\sigma>0\), define the endpoint regularized power chain operator
\begin{equation}\label{eq:endpoint-regularized-channel-operator}
\boxed{\;
(W_{\mathrm{pow}}^{(\sigma)}g)(0)
:=\sum_{d\ge 1}\frac{1}{d^{1+\sigma}}\,g\!\bigl(C_d(0)\bigr)\; }.
\end{equation}
Then its endpoint three-channel decomposition is
\begin{equation}\label{eq:endpoint-regularized-channel-decomposition}
\boxed{\;
(W_{\mathrm{pow}}^{(\sigma)}g)(0)
=\mathsf H_{\mathrm{odd}}(\sigma)\,g(0)
+\mathsf H_{0}(\sigma)\,g(2)
+\mathsf H_{2}(\sigma)\,g(-2)\; }.
\end{equation}
The three-channel weight kernels admit the following closed-form Dirichlet symbol representations:
\begin{align}
\mathsf H_{\mathrm{odd}}(\sigma)
&:=\sum_{\substack{d\ge 1\\ d\ \mathrm{odd}}}\frac{1}{d^{1+\sigma}}
=\bigl(1-2^{-(1+\sigma)}\bigr)\zeta(1+\sigma),\label{eq:endpoint-Hodd-sigma}\\
\mathsf H_{0}(\sigma)
&:=\sum_{d\equiv 0\ (\mathrm{mod}\ 4)}\frac{1}{d^{1+\sigma}}
=4^{-(1+\sigma)}\zeta(1+\sigma),\label{eq:endpoint-H0-sigma}\\
\mathsf H_{2}(\sigma)
&:=\sum_{d\equiv 2\ (\mathrm{mod}\ 4)}\frac{1}{d^{1+\sigma}}
=\bigl(2^{-(1+\sigma)}-4^{-(1+\sigma)}\bigr)\zeta(1+\sigma).\label{eq:endpoint-H2-sigma}
\end{align}
\end{proposition}
\begin{proof}
By the endpoint tristate compression \eqref{eq:endpoint-tristate-values}, \(C_d(0)\) depends only on \(d\bmod 4\) and takes values in \(\{0,2,-2\}\), hence
\[
(W_{\mathrm{pow}}^{(\sigma)}g)(0)
=\sum_{\substack{d\ge 1\\ d\ \mathrm{odd}}}\frac{1}{d^{1+\sigma}}\,g(0)
+\sum_{d\equiv 0(4)}\frac{1}{d^{1+\sigma}}\,g(2)
+\sum_{d\equiv 2(4)}\frac{1}{d^{1+\sigma}}\,g(-2),
\]
yielding \eqref{eq:endpoint-regularized-channel-decomposition}.
For the closed-form representations, note that
\[
\sum_{\substack{d\ge 1\\ d\ \mathrm{odd}}}\frac{1}{d^{1+\sigma}}
\;=\;\sum_{d\ge 1}\frac{1}{d^{1+\sigma}}-\sum_{k\ge 1}\frac{1}{(2k)^{1+\sigma}}
\;=\;\bigl(1-2^{-(1+\sigma)}\bigr)\zeta(1+\sigma),
\]
and
\[
\sum_{d\equiv 0(4)}\frac{1}{d^{1+\sigma}}
=\sum_{k\ge 1}\frac{1}{(4k)^{1+\sigma}}
=4^{-(1+\sigma)}\zeta(1+\sigma),
\]
as well as
\(
\sum_{d\equiv 2(4)}d^{-(1+\sigma)}
=\sum_{k\ge 1}(2k)^{-(1+\sigma)}-\sum_{k\ge 1}(4k)^{-(1+\sigma)}
\)
which gives \eqref{eq:endpoint-H2-sigma}. \qed
\end{proof}

\paragraph{Four-Channel Lifting of the Endpoint Phase: \texorpdfstring{$\ZZ/4$}{Z/4} Branch Memory and Its Minimal Closure in the Completion Layer}
In the completion coordinates \(u=r^{2},\ S=r+r^{-1}\), the invariant subring under the inversion \(r\mapsto r^{-1}\) descends strictly to \(\ZZ[S]\) (Proposition \ref{prop:completion-subring-unique-angle}), and therefore the visible information at the endpoint \(u=-1\) degenerates at the \(0\)-jet level to the tristate \(\mathcal S_\star=\{0,\pm 2\}\) (Corollary \ref{cor:endpoint-tristate-compression}).
To capture the \(\ZZ/4\) branch memory at the endpoint within the completion layer, we introduce the minimal inversion-anti-invariant adjoint coordinate
\begin{equation}\label{eq:endpoint-delta-anti-invariant}
\boxed{\ \delta:=r-r^{-1}\ }.
\end{equation}
and for any \(d\ge 1\) define the "odd" Chebyshev--Adams adjoint polynomial family
\begin{equation}\label{eq:endpoint-odd-chebyshev-adams}
\boxed{\ D_d(S,\delta):=r^{d}-r^{-d}\ }.
\end{equation}

\begin{remark}[Principal Generator of the Inversion-Anti-Invariant Sector and Minimal Closure]\label{rem:endpoint-inv-antiinv-min-closure}
Let
\[
\mathcal A^{-}:=\{\,f(r)\in \ZZ[r,r^{-1}]:\ f(r^{-1})=-f(r)\,\}
\]
be the anti-invariant part under the inversion action \(r\mapsto r^{-1}\).
Then
\[
\boxed{\ \mathcal A^{-}=\delta\,\ZZ[S]\ },\qquad S=r+r^{-1},\ \delta=r-r^{-1}.
\]
Therefore, at the completion level, adjoining \(\delta\) as an additional fiber coordinate generates all anti-invariant observables, and at the endpoint branches \(r=\pm\mathrm{i}\) it reads out the branch memory as \(\delta=\pm 2\mathrm{i}\); this provides the minimal closure from the tristate \(\mathcal S_\star\) to the complete \(\ZZ/4\) channel system.
\end{remark}

\begin{proposition}[Odd Chebyshev--Adams Adjoint Identity]\label{prop:endpoint-odd-chebyshev-identity}
For any \(d\ge 1\), the identity
\begin{equation}\label{eq:endpoint-odd-chebyshev-identity}
\boxed{\;
D_d(S,\delta)=\delta\,U_{d-1}(S/2)\; },
\end{equation}
holds, where \(U_{d-1}\) is the Chebyshev polynomial of the second kind.
In particular, at the endpoint \(u=-1\) on the half-angle branches \(r=\pm\mathrm{i}\), we have \(S=0\), \(\delta=\pm 2\mathrm{i}\), and
\begin{equation}\label{eq:endpoint-odd-chebyshev-ratio}
\boxed{\;
\frac{D_d(0,\delta)}{\delta}
=\frac{r^{d}-r^{-d}}{r-r^{-1}}
=U_{d-1}(0)
=\sin(d\pi/2)\; }.
\end{equation}
\end{proposition}
\begin{proof}
Let \(E_d:=r^{d}-r^{-d}\). Then from \(r+r^{-1}=S\) one directly verifies the recurrence
\(
E_{d+1}=SE_d-E_{d-1}
\)
with initial values \(E_0=0,E_1=\delta\).
Therefore \(E_d/\delta\) satisfies the same recurrence and initial conditions as \(U_{d-1}(S/2)\) (\(U_0=1,U_1=S\)), yielding \eqref{eq:endpoint-odd-chebyshev-identity}.
\par
Setting \(S=0\) and using the identity \(U_{d-1}(\cos\phi)=\sin(d\phi)/\sin\phi\) (with \(\phi=\pi/2\)) gives \(U_{d-1}(0)=\sin(d\pi/2)\), yielding \eqref{eq:endpoint-odd-chebyshev-ratio}. \qed
\end{proof}

\begin{theorem}[Endpoint Four-Channel Closure and \texorpdfstring{$\chi_{4}$}{chi4} Readout]\label{thm:endpoint-four-channel-chi4}
Define the endpoint channel index
\begin{equation}\label{eq:endpoint-epsilon-def}
\boxed{\ \varepsilon(d):=\sin(d\pi/2)\in\{-1,0,1\}\qquad(d\ge 1)\ }.
\end{equation}
Then \(\varepsilon(d)=0\) if and only if \(d\) is even; when \(d\) is odd,
\[
\varepsilon(d)=
\begin{cases}
1, & d\equiv 1\ (\mathrm{mod}\ 4),\\
-1,& d\equiv 3\ (\mathrm{mod}\ 4).
\end{cases}
\]
Therefore \(\varepsilon\) provides the complete odd-class splitting modulo \(4\).
Let \(\chi_{4}\) be the nontrivial Dirichlet character modulo \(4\):
\[
\boxed{\;
\chi_{4}(d)=
\begin{cases}
0, & d\ \text{is even},\\
1, & d\equiv 1\ (\mathrm{mod}\ 4),\\
-1,& d\equiv 3\ (\mathrm{mod}\ 4),
\end{cases}\;}
\]
Then for all \(d\ge 1\) we have \(\chi_{4}(d)=\varepsilon(d)\).
This result shows that the completion even part \(C_d\) alone yields only the endpoint tristate, whereas adjoining the minimal inversion-anti-invariant fiber coordinate \(\delta\) naturally closes the endpoint phase branch memory to the complete \(\ZZ/4\) channel system.
\end{theorem}
\begin{proof}
Since the values of \(\sin(d\pi/2)\) depend only on \(d\bmod 4\), direct enumeration of the four residue classes yields the stated splitting; comparison with the class-by-class definition of \(\chi_4\) shows \(\chi_4(d)=\varepsilon(d)\). \qed
\end{proof}

\paragraph{Finite-Dimensional Invariant Subrepresentations at the Endpoint: The Binary Symbol of \texorpdfstring{$\zeta$}{zeta} and \texorpdfstring{$L(\chi_{4})$}{L(chi4)}}
For \(\sigma>0\), introduce the regularized Chebyshev--Witt power chain operator in completion coordinates
\begin{equation}\label{eq:endpoint-Wsigma-even}
\boxed{\;
(\mathcal W_{\sigma}F)(S):=\sum_{m\ge 1}\frac{1}{m^{1+\sigma}}\,F\!\bigl(C_m(S)\bigr)\; },
\qquad S\in[-2,2],\ F\in\CC[S],
\end{equation}
and its adjoint action on the anti-invariant fiber extension \((S,\delta)\):
\begin{equation}\label{eq:endpoint-Wsigma-odd}
\boxed{\;
(\mathcal W_{\sigma}^{\delta}G)(S,\delta):=\sum_{m\ge 1}\frac{1}{m^{1+\sigma}}\,G\!\bigl(C_m(S),D_m(S,\delta)\bigr)\; },
\qquad S\in[-2,2],\ G\in\CC[S,\delta].
\end{equation}
Define three even endpoint linear functionals and one anti-invariant endpoint linear functional:
\[
\mathcal E_{0}(F):=F(0),\qquad \mathcal E_{+}(F):=F(2),\qquad \mathcal E_{-}(F):=F(-2),
\]
\begin{equation}\label{eq:endpoint-odd-functional-O}
\boxed{\ \mathcal O(G):=\frac{1}{2\mathrm{i}}\Bigl(G(0,2\mathrm{i})-G(0,-2\mathrm{i})\Bigr)\ }.
\end{equation}

\begin{theorem}[Endpoint Symbol Decomposition: The Orthogonal Pair \texorpdfstring{$\zeta$}{zeta}--\texorpdfstring{$L(\chi_{4})$}{L(chi4)}]\label{thm:endpoint-zeta-Lchi4-symbol-pair}
For any \(\sigma>0\) and \(F\in\CC[S]\), the endpoint even observable decomposes as
\begin{equation}\label{eq:endpoint-even-symbol-decomposition-alpha}
\boxed{\;
\mathcal E_{0}(\mathcal W_{\sigma}F)
=\alpha_{0}(\sigma)\,\mathcal E_{0}(F)+\alpha_{+}(\sigma)\,\mathcal E_{+}(F)+\alpha_{-}(\sigma)\,\mathcal E_{-}(F)\; },
\end{equation}
where
\[
\alpha_{0}(\sigma):=\sum_{\substack{m\ge 1\\ m\ \mathrm{odd}}}\frac{1}{m^{1+\sigma}},\qquad
\alpha_{+}(\sigma):=\sum_{m\equiv 0(4)}\frac{1}{m^{1+\sigma}},\qquad
\alpha_{-}(\sigma):=\sum_{m\equiv 2(4)}\frac{1}{m^{1+\sigma}}.
\]
Moreover, \(\alpha_{0},\alpha_{+},\alpha_{-}\) admit closed-form \(\zeta\)-symbol representations (Proposition \ref{prop:endpoint-dirichlet-symbol-weights}, equations \eqref{eq:endpoint-Hodd-sigma}--\eqref{eq:endpoint-H2-sigma}).
\par
Simultaneously, for any \(G\in\CC[S,\delta]\), the endpoint anti-invariant observable satisfies the eigenvalue decomposition
\begin{equation}\label{eq:endpoint-odd-symbol-eigen}
\boxed{\;
\mathcal O(\mathcal W_{\sigma}^{\delta}G)
=L(1+\sigma,\chi_{4})\cdot \mathcal O(G)\; },
\end{equation}
where
\(
L(1+\sigma,\chi_{4}):=\sum_{m\ge 1}\chi_{4}(m)m^{-(1+\sigma)}
\)
is the Dirichlet \(L\)-series of the modulo \(4\) character.
In other words, the endpoint anti-invariant observable \(\mathcal O\) constitutes a one-dimensional invariant subrepresentation under \(\mathcal W_{\sigma}^{\delta}\), whose symbol is precisely \(L(1+\sigma,\chi_{4})\); it is orthogonally decoupled from the \(\zeta\)-type weight kernels of the even channels.
\end{theorem}
\begin{proof}
The even observable \eqref{eq:endpoint-even-symbol-decomposition-alpha} is obtained by classifying the summation according to \(m\bmod 4\) via the endpoint tristate compression \eqref{eq:endpoint-tristate-values}; the closed-form expressions are from Proposition \ref{prop:endpoint-dirichlet-symbol-weights}.
\par
For the anti-invariant observable, using \eqref{eq:endpoint-odd-functional-O} and \eqref{eq:endpoint-Wsigma-odd} gives
\[
\mathcal O(\mathcal W_{\sigma}^{\delta}G)
=\frac{1}{2\mathrm{i}}\sum_{m\ge 1}\frac{1}{m^{1+\sigma}}
\Bigl(G(C_m(0),D_m(0,2\mathrm{i}))-G(C_m(0),D_m(0,-2\mathrm{i}))\Bigr).
\]
If \(m\) is even, then by \eqref{eq:endpoint-epsilon-def} we have \(\varepsilon(m)=0\), and by Proposition \ref{prop:endpoint-odd-chebyshev-identity}, \(D_m(0,\pm 2\mathrm{i})=0\), so the corresponding terms cancel exactly in the difference.
If \(m\) is odd, then \(C_m(0)=0\) and \(\varepsilon(m)=\pm 1\), and by \eqref{eq:endpoint-odd-chebyshev-ratio},
\(
D_m(0,\pm 2\mathrm{i})=\pm 2\mathrm{i}\,\varepsilon(m)
\).
Hence
\[
G(0,2\mathrm{i}\varepsilon(m))-G(0,-2\mathrm{i}\varepsilon(m))
=2\mathrm{i}\,\varepsilon(m)\,\mathcal O(G).
\]
Substituting back yields
\[
\mathcal O(\mathcal W_{\sigma}^{\delta}G)
=\left(\sum_{m\ge 1}\frac{\varepsilon(m)}{m^{1+\sigma}}\right)\mathcal O(G)
=\left(\sum_{m\ge 1}\frac{\chi_{4}(m)}{m^{1+\sigma}}\right)\mathcal O(G)
=L(1+\sigma,\chi_{4})\,\mathcal O(G),
\]
where the second equality follows from Theorem \ref{thm:endpoint-four-channel-chi4} (\(\varepsilon=\chi_4\)). \qed
\end{proof}

\begin{proposition}[Invariant Subrepresentation and Explicit Matrix Form of the Even Tristate at the Endpoint]\label{prop:endpoint-even-tristate-matrix}
Using the notation from \eqref{eq:endpoint-Wsigma-even}--\eqref{eq:endpoint-odd-symbol-eigen}, let
\[
v(F):=\bigl(\mathcal E_{0}(F),\mathcal E_{+}(F),\mathcal E_{-}(F)\bigr)^{\mathsf T}\in\CC^{3}.
\]
Then for any \(\sigma>0\) and any \(F\in\CC[S]\), the endpoint even tristate is closed under \(\mathcal W_\sigma\), satisfying
\begin{equation}\label{eq:endpoint-even-tristate-matrix-action}
\boxed{\;
v(\mathcal W_{\sigma}F)=M_{\mathrm{even}}(\sigma)\,v(F)\;}
\end{equation}
where the representation matrix has the closed form
\begin{equation}\label{eq:endpoint-even-tristate-matrix-closed}
\boxed{\;
M_{\mathrm{even}}(\sigma)
=\zeta(1+\sigma)
\begin{pmatrix}
1-2^{-(1+\sigma)} & 4^{-(1+\sigma)} & 2^{-(1+\sigma)}-4^{-(1+\sigma)}\\
0&1&0\\
0&2^{-(1+\sigma)}&1-2^{-(1+\sigma)}
\end{pmatrix}. \;}
\end{equation}
\end{proposition}

\begin{proof}
The first row is given by the endpoint even observable decomposition \eqref{eq:endpoint-even-symbol-decomposition-alpha} of Theorem \ref{thm:endpoint-zeta-Lchi4-symbol-pair} together with the closed-form expressions of the \(\alpha\)-weight kernels.
\par
For \(\mathcal E_{+}\) and \(\mathcal E_{-}\), the doubling-angle identities of the completion trace coordinate yield \(C_m(2)=2\) and \(C_m(-2)=2(-1)^m\), so
\[
\mathcal E_{+}(\mathcal W_\sigma F)
=\sum_{m\ge 1}\frac{1}{m^{1+\sigma}}\,F\!\bigl(C_m(2)\bigr)
=\zeta(1+\sigma)\,\mathcal E_{+}(F),
\]
\[
\mathcal E_{-}(\mathcal W_\sigma F)
=\sum_{\substack{m\ge 1\\ m\ \mathrm{even}}}\frac{1}{m^{1+\sigma}}\,F(2)
\;+\;
\sum_{\substack{m\ge 1\\ m\ \mathrm{odd}}}\frac{1}{m^{1+\sigma}}\,F(-2)
=2^{-(1+\sigma)}\zeta(1+\sigma)\,\mathcal E_{+}(F)
\;+\;\bigl(1-2^{-(1+\sigma)}\bigr)\zeta(1+\sigma)\,\mathcal E_{-}(F).
\]
Writing the above three equations in matrix form yields \eqref{eq:endpoint-even-tristate-matrix-action}--\eqref{eq:endpoint-even-tristate-matrix-closed}. \qed
\end{proof}

\begin{theorem}[Universal Pole and Finite Part of the Even Endpoint Channel (\texorpdfstring{$\sigma\downarrow 0$}{sigma->0})]\label{thm:endpoint-even-finite-part-renormalization}
Let
\[
\mu:=\left(\frac12,\frac14,\frac14\right),\qquad
\langle \mu,(v_0,v_+,v_-)\rangle:=\frac12 v_0+\frac14 v_++\frac14 v_-.
\]
For any \(F\in\CC[S]\), define the renormalized finite part functional of the even endpoint channel
\begin{equation}\label{eq:endpoint-even-renormalized-finite-part-def}
\boxed{\;
\mathcal R_{\mathrm{even}}(F)
:=\lim_{\sigma\downarrow 0}\left(
\mathcal E_{0}(\mathcal W_\sigma F)
-\frac{1}{\sigma}\,\langle \mu,v(F)\rangle
\right)\;}
\end{equation}
(provided the limit exists).
Then this limit exists and yields the explicit finite part
\begin{equation}\label{eq:endpoint-even-renormalized-finite-part-closed}
\boxed{\;
\mathcal R_{\mathrm{even}}(F)
=\left(\frac{\gamma}{2}+\frac{\log 2}{2}\right)\mathcal E_{0}(F)
\;+\;\left(\frac{\gamma}{4}-\frac{\log 2}{2}\right)\mathcal E_{+}(F)
\;+\;\frac{\gamma}{4}\mathcal E_{-}(F)\; }.
\end{equation}
where \(\gamma\) is the Euler constant.
\end{theorem}

\begin{proof}
By Theorem \ref{thm:endpoint-zeta-Lchi4-symbol-pair} and \eqref{eq:endpoint-even-symbol-decomposition-alpha},
\[
\mathcal E_0(\mathcal W_\sigma F)
=\alpha_0(\sigma)\mathcal E_0(F)+\alpha_+(\sigma)\mathcal E_+(F)+\alpha_-(\sigma)\mathcal E_-(F),
\]
and (Proposition \ref{prop:endpoint-dirichlet-symbol-weights})
\[
\alpha_0(\sigma)=\bigl(1-2^{-(1+\sigma)}\bigr)\zeta(1+\sigma),\quad
\alpha_+(\sigma)=4^{-(1+\sigma)}\zeta(1+\sigma),\quad
\alpha_-(\sigma)=\bigl(2^{-(1+\sigma)}-4^{-(1+\sigma)}\bigr)\zeta(1+\sigma).
\]
Using the Laurent expansion \(\zeta(1+\sigma)=\sigma^{-1}+\gamma+O(\sigma)\) (\(\sigma\downarrow 0\)) together with
\[
2^{-(1+\sigma)}=\frac12\bigl(1-\sigma\log 2+O(\sigma^2)\bigr),\qquad
4^{-(1+\sigma)}=\frac14\bigl(1-2\sigma\log 2+O(\sigma^2)\bigr),
\]
we obtain
\[
\alpha_0(\sigma)=\frac{1}{2\sigma}+\left(\frac{\gamma}{2}+\frac{\log 2}{2}\right)+O(\sigma),\quad
\alpha_+(\sigma)=\frac{1}{4\sigma}+\left(\frac{\gamma}{4}-\frac{\log 2}{2}\right)+O(\sigma),\quad
\alpha_-(\sigma)=\frac{1}{4\sigma}+\frac{\gamma}{4}+O(\sigma).
\]
Substituting back and separating the \(\sigma^{-1}\) polar term from the constant term yields the existence of the limit \eqref{eq:endpoint-even-renormalized-finite-part-def} equal to \eqref{eq:endpoint-even-renormalized-finite-part-closed}. \qed
\end{proof}

\begin{corollary}[Analyticity of the Anti-Invariant Endpoint Channel and \texorpdfstring{$\chi_4$}{chi4} Constant Output]\label{cor:endpoint-odd-analytic-at-zero}
For any \(G\in\CC[S,\delta]\),
\begin{equation}\label{eq:endpoint-odd-limit-L1-chi4}
\boxed{\;
\lim_{\sigma\downarrow 0}\mathcal O(\mathcal W_{\sigma}^{\delta}G)
=L(1,\chi_{4})\cdot \mathcal O(G)
=\frac{\pi}{4}\,\mathcal O(G)\; }.
\end{equation}
\end{corollary}

\begin{proof}
By the eigenvalue decomposition \eqref{eq:endpoint-odd-symbol-eigen} of Theorem \ref{thm:endpoint-zeta-Lchi4-symbol-pair}, \(\mathcal O(\mathcal W_{\sigma}^{\delta}G)=L(1+\sigma,\chi_4)\mathcal O(G)\).
Since \(L(1+\sigma,\chi_4)\) is analytic at \(\sigma=0\), taking the limit yields the first equality; the classical identity \(L(1,\chi_4)=\pi/4\) gives the second equality. \qed
\end{proof}

\begin{remark}[Binary Residue of the Endpoint Renormalization: Universal Pole \(\oplus\) Discrete Eigenvalue Channel]\label{rem:endpoint-renormalization-universal-pole-discrete-hair}
Theorem \ref{thm:endpoint-even-finite-part-renormalization} shows that all divergences of the even endpoint channel occur along the rank-one direction
\(\langle\mu,v(F)\rangle\) and are entirely controlled by the pole of \(\zeta(1+\sigma)\) at \(1\);
the finite part is given by an explicit combination of the Euler constant and \(\log 2\).
Orthogonally, Corollary \ref{cor:endpoint-odd-analytic-at-zero} shows that the anti-invariant endpoint observable has no pole at \(\sigma\downarrow 0\), and its residue compresses to a purely discrete \(\chi_4\) constant channel \(L(1,\chi_4)=\pi/4\).
Therefore, under the endpoint symbol caliber, the non-compactness after renormalization can be compressed into a binary orthogonal decomposition: "universal \(\zeta\)-pole + discrete \(\chi_4\) eigenvalue residue."
\end{remark}

\begin{remark}[Endpoint Logarithmic Principal Scale and Phase Precision Decoupling: Counting Control of the Even Power Channel at \texorpdfstring{$u=-1$}{u=-1}]\label{rem:endpoint-logscale-parity-decouple}
On the twist layer \(u=-1\), even powers satisfy \(u^{2k}=1\), so any Witt-type dilation superposition
\(
\sum_{m\ge 1}\frac1m f(u^m)
\)
contains a logarithmic principal scale induced by the even-power subchain \(\sum_{k\ge 1}\frac{1}{2k}f(1)\).
For primitive input \(f=p_n\), the coefficient of this principal scale is
\(
f(1)=p_n(1)=B_n^{+}
\)
(see \eqref{eq:primitive-endpoint-Bpm-identification}), which is by definition separated from the half-angle branch information encoded by the endpoint difference certificate
\(
A_n=\widehat p_n(0)
\).
In particular, when \(n\) is odd, \(\widehat p_n\) is an odd polynomial (Corollary \ref{cor:primitive-completion-hatp}), so \(A_n=\widehat p_n(0)=0\); and when \(n\ge 3\) is odd, we also have \(p_n(-1)=0\) (Corollary \ref{cor:primitive-odd-vanish-u-minus1}), so the endpoint parity difference certificate vanishes completely on odd-length primitives.
\end{remark}

\paragraph{Frobenius--Adams Integer-Layer Phase Certificate: Stratified Audit of \texorpdfstring{$\Delta_{p,k}$}{Delta} on the Endpoint Tristate}
Definition \ref{def:witt-frobenius-defect} gives the Frobenius defect polynomial
\(
\Delta_{p,k}(u):=a_{p^k}(u)-a_{p^{k-1}}(u^p)\in\ZZ[u]
\),
and Proposition \ref{prop:witt-frobenius-defect-divisibility} yields the divisibility depth
\(
\Delta_{p,k}(u)\in p^k\ZZ[u]
\).
On the other hand, the coefficient ring version of the Big-Witt/M\"obius inversion on a commutative ring \(R\) equipped with Adams operations \(\psi^d\) satisfies
\(
n\,p_n=\sum_{d\mid n}\mu(d)\,\psi^d(a_{n/d})\in R
\)
(Proposition \ref{prop:witt-necklace-dwork-R}), where in the multi-parameter polynomial ring case \(\psi^d\) is realized by variable-wise power substitution.
Juxtaposing the above two facts with the endpoint tristate compression, one can define a family of integer certificates whose realization chain is directly auditable.

\begin{definition}[Frobenius Integer Certificate on the Endpoint Tristate]\label{def:endpoint-frobenius-integer-certificate}
For any prime \(p\), integer \(k\ge 1\), and endpoint state \(\sigma\in\mathcal S_{\star}\), choose \(r_\sigma\in\{\,\mathrm{i},1,-1\,\}\) such that \(r_\sigma+r_\sigma^{-1}=\sigma\), and set \(u_\sigma:=r_\sigma^{2}\in\{-1,1\}\).
Define
\begin{equation}\label{eq:endpoint-frobenius-integer-certificate}
\boxed{\ \mathfrak I_{p,k}(\sigma):=\frac{\Delta_{p,k}(u_{\sigma})}{p^{k}}\in\ZZ\ }.
\end{equation}
\end{definition}

\begin{proposition}[Integer-Layer A Priori Guardrail: Auditability of the Endpoint Certificate]\label{prop:endpoint-frobenius-integer-certificate-auditable}
For any \((p,k)\) and \(\sigma\in\mathcal S_{\star}\), the certificate \(\mathfrak I_{p,k}(\sigma)\) is a well-defined integer.
If in any realization \(\mathfrak I_{p,k}(\sigma)\notin\ZZ\) (equivalently \(\Delta_{p,k}(u_\sigma)\notin p^k\ZZ\)), then the Frobenius defect divisibility depth \(\Delta_{p,k}(u)\in p^k\ZZ[u]\) is necessarily violated, so that errors can be localized at the integer layer before entering the discussion of continuous spectral layers such as the unit circle spectral radius.
\end{proposition}
\begin{proof}
By Proposition \ref{prop:witt-frobenius-defect-divisibility}, all coefficients of \(\Delta_{p,k}(u)\) are divisible by \(p^k\); therefore for any \(u_\sigma\in\ZZ\) we have \(\Delta_{p,k}(u_\sigma)\in p^k\ZZ\), hence \(\mathfrak I_{p,k}(\sigma)\in\ZZ\).
If a counterexample arises, it contradicts "coefficient-level divisibility." \qed
\end{proof}

\paragraph{Orthogonality of Integer-Layer Residuals and Spectral-Layer Shifts: Frobenius Defect as an Arithmetic Precision Axis}
From \(\Delta_{p,k}(u)\in p^k\ZZ[u]\), one can define the normalized residual polynomial at the integer layer:
\begin{equation}\label{eq:frobenius-normalized-residual-poly}
\boxed{\ E_{p,k}(u):=\frac{\Delta_{p,k}(u)}{p^{k}}\in\ZZ[u]\ }.
\end{equation}
The family \(\{E_{p,k}\}_{p,k}\) thus constitutes a purely discrete "arithmetic precision axis": its audit depends only on integer coefficients and divisibility depth, and is independent of continuous spectral layer quantities such as the unit circle spectral radius and endpoint precision potential.
On the endpoint tristate \(\sigma\in\mathcal S_{\star}\), the above residuals are strictly consistent with the certificates of Definition \ref{def:endpoint-frobenius-integer-certificate}:
\[
\boxed{\ \mathfrak I_{p,k}(\sigma)=E_{p,k}(u_\sigma)\in\ZZ\ }.
\]
Consequently, the principle of "integer layer first, spectral layer second" in stratified auditing can be formulated as follows: first lock down the consistency guardrail of a realization via \eqref{eq:frobenius-normalized-residual-poly} (or its endpoint evaluations),
then enter the analytical layer discussion involving the endpoint precision potential \(y=-\log|\mathfrak{w}|\) (equation \eqref{eq:endpoint-busemann-potential}) and continuous spectral shifts on the unitary slice.
